\documentclass[11pt]{article}
\usepackage[margin=1in]{geometry}
\usepackage[T1]{fontenc}
\usepackage{lmodern}
\usepackage{amsmath,amssymb,amsthm}
\usepackage{graphicx}
\usepackage[hidelinks]{hyperref}

\newtheorem{theorem}{Theorem}
\newcommand{\D}{\mathbb D}
\newcommand{\T}{\mathbb T}
\newcommand{\N}{\mathbb N}
\newcommand{\cK}{K}

\title{A Literature-Implied Positive Answer to Question 2.5 of
arXiv:2102.03447:\\Model-Space Bessel Systems as a Vectorial Hankel Operator}
\author{Run fa\_banach\_001, agent\_lane\_05}
\date{11 August 2026}

\begin{document}
\maketitle

\section*{Status and source question}

This is a \texttt{literature\_implied\_answer (full positive answer)} packet.
Dayan asks in Question~2.5, on PDF page~8 of \cite{Dayan}, whether
\begin{equation}\label{eq:condition}
 C:=\sup_{z\in\D}\sum_{n\in\N}\bigl(1-|\Theta_n(z)|^2\bigr)<\infty
\end{equation}
forces the model spaces
\(\cK_{\Theta_n}=H^2\ominus\Theta_nH^2\) to form a Bessel system, and asks
separately whether this holds when all \(\Theta_n\) are Blaschke products.

\begin{center}
\includegraphics[width=.96\textwidth]{figures/open_problem_crop.png}
\end{center}

The answer is yes for arbitrary inner functions.  The supporting result is
Treil's reproducing-kernel thesis for possibly vectorial pre-Hankel operators
\cite[Theorem~1.1]{Treil}.  Treil's theorem predates the source question and
does not mention model spaces; the implication below is the direct
operator-theoretic identification made in this run.

\section*{Identification}

Let \(P_-:L^2(\T)\to H^2_-=L^2(\T)\ominus H^2\) denote projection onto the
strictly negative Fourier frequencies, and put \(E=\ell^2(\N)\).  For an inner
function \(\Theta\),
\begin{equation}\label{eq:projection}
 \Theta P_-(\overline\Theta f)
 =f-\Theta P_+(\overline\Theta f)
 =P_{\cK_\Theta}f,
 \qquad f\in H^2.
\end{equation}
Multiplication by \(\Theta\) is isometric, so
\begin{equation}\label{eq:normidentity}
 \|P_-(\overline\Theta f)\|_2=\|P_{\cK_\Theta}f\|_2.
\end{equation}

The hypothesis at \(z=0\) makes the vectorial construction legitimate.  In
fact,
\[
 P_-\overline{\Theta_n}
 =\overline{\Theta_n}-\overline{\Theta_n(0)},
 \qquad
 \|P_-\overline{\Theta_n}\|_2^2=1-|\Theta_n(0)|^2,
\]
and hence
\[
 \Phi:=\bigl(P_-\overline{\Theta_n}\bigr)_{n\in\N}
 \in H^2_-(E).
\]
It is therefore the antianalytic symbol of the vectorial pre-Hankel operator
\begin{equation}\label{eq:Gamma}
 \Gamma f=P_-(\Phi f)
 =\bigl(P_-(\overline{\Theta_n}f)\bigr)_{n\in\N},
\end{equation}
initially defined on analytic polynomials.  The second equality holds because
subtracting the constant \(\overline{\Theta_n(0)}\) does not affect \(P_-\)
after multiplication by an analytic \(f\).

\section*{Full answer}

\begin{theorem}
Let \((\Theta_n)_{n\in\N}\) be any sequence of inner functions satisfying
\eqref{eq:condition}.  Then \((\cK_{\Theta_n})_{n\in\N}\) is a Bessel system,
and one may take Bessel bound \(4eC\):
\[
 \sum_{n\in\N}\|P_{\cK_{\Theta_n}}f\|_2^2
 \le 4eC\,\|f\|_2^2,
 \qquad f\in H^2.
\]
Conversely, a Bessel bound \(B\) implies \eqref{eq:condition} with \(C\le B\).
Thus the kernel condition characterizes Bessel systems of model spaces.
\end{theorem}

\begin{proof}
Let
\[
 k_z(w)=\frac{(1-|z|^2)^{1/2}}{1-\overline zw}
\]
be the normalized Szeg\H{o} kernel.  For an inner \(\Theta\), the adjoint
multiplier identity \(M_\Theta^*k_z=\overline{\Theta(z)}k_z\) gives
\begin{equation}\label{eq:kernelnorm}
 \|P_{\cK_\Theta}k_z\|_2^2
 =\langle P_{\cK_\Theta}k_z,k_z\rangle
 =1-|\Theta(z)|^2.
\end{equation}
Combining \eqref{eq:normidentity}, \eqref{eq:Gamma}, and
\eqref{eq:kernelnorm},
\[
 \|\Gamma k_z\|_{H^2_-(E)}^2
 =\sum_n\|P_{\cK_{\Theta_n}}k_z\|_2^2
 =\sum_n\bigl(1-|\Theta_n(z)|^2\bigr)
 \le C.
\]

Treil's theorem applies because the domain is scalar \(H^2\), the target is
the vector-valued space \(H^2_-(E)\), and \(\Gamma\) is a vectorial
pre-Hankel operator.  It yields
\[
 \|\Gamma\|\le 2\sqrt e\,\sqrt C.
\]
For analytic polynomials \(f\), equations \eqref{eq:normidentity} and
\eqref{eq:Gamma} now give
\[
 \sum_n\|P_{\cK_{\Theta_n}}f\|_2^2
 =\|\Gamma f\|_{H^2_-(E)}^2
 \le4eC\|f\|_2^2.
\]
Density extends the inequality to every \(f\in H^2\).  Conversely, applying a
Bessel inequality to the unit vector \(k_z\) and using
\eqref{eq:kernelnorm} gives the displayed kernel condition, uniformly in
\(z\).  This proves both assertions.
\end{proof}

\begin{center}
\includegraphics[width=.94\textwidth]{figures/supporting_theorem_crop.png}
\end{center}

\section*{Provenance, search bounds, and limitations}

Treil's arXiv revision explicitly says that its proof covers vectorial Hankel
operators, and Theorem~1.1 states the needed scalar-domain/vector-target
result with the constant \(2\sqrt e\).  Dayan's bibliography does not cite
Treil, and the supporting paper predates and does not explicitly discuss
Dayan's question.  Accordingly this packet records an agent-identified
literature implication, not a newly proved analytic theorem.

The run's cheap indexes were searched for the source arXiv id, title, and core
keywords.  A bounded web/arXiv search used the exact displayed condition and
question wording, together with ``model-space Bessel system,'' ``vectorial
Hankel operator,'' and ``reproducing kernel thesis.''  No later paper explicitly
answering Question~2.5 was found; the reformulation led to \cite{Treil}.

There is no Blaschke-product limitation: the proof covers every countable
sequence of inner functions.  The inherited constant \(4e\) is explicit but
is not asserted to be optimal.  Human review should focus on
\eqref{eq:projection} and the vectorial-symbol construction; these are the only
identifications beyond the cited theorem.

\begin{thebibliography}{9}

\bibitem{Dayan}
A.~Dayan,
\emph{Weakly Separated Bessel Systems of Model Spaces},
arXiv:2102.03447v2 (2021), Question~2.5,
\href{https://arxiv.org/abs/2102.03447}{arXiv:2102.03447}.

\bibitem{Treil}
S.~Treil,
\emph{A remark on the reproducing kernel thesis for Hankel operators},
arXiv:1201.0063v2 (2012), Theorem~1.1;
St. Petersburg Math. J. \textbf{26} (2015), 479--485,
\href{https://arxiv.org/abs/1201.0063}{arXiv:1201.0063}.

\end{thebibliography}

\end{document}

